% EDITORIAL QUESTION: Como evoluíram historicamente as escalas de graduação na Europa, EUA e Brasil até os polos modernos de referência?
\section{Evolução Histórico-Geográfica das Escalas}

Os sistemas atuais resultam de um século de transição entre escalas fechadas e abertas.

\subsection*{A.1 Das Escalas Alpinas Fechadas à Abertura de Yosemite e da UIAA}
% SOURCE [VERIFIED_PRIMARY]: UIAA History (https://www.theuiaa.org/inside-the-uiaa/history/); Sportrock (https://www.sportrock.com/post/understanding-climbing-grades).
Em \textbf{1926}, Willo Welzenbach propôs nos Alpes uma escala fechada de Grau I a Grau VI \cite{uiaa_history}. Em \textbf{1935}, o Sierra Club criou o \textbf{Sistema Decimal de Yosemite (YDS)}, dividindo a montanha em cinco classes e prevendo a Classe 5 fechada entre \texttt{5.0} e \texttt{5.9}. A evolução técnica em Yosemite forçou a abertura para \texttt{5.10a--d} até \texttt{5.15d}, e a UIAA aboliu o teto do Grau VI em 1978 \cite{uiaa_history, sportrock_grades}.

\subsection*{A.2 Fontainebleau: Circuitos por Cores e Escala Font}
% SOURCE [VERIFIED_PRIMARY]: Climb Ski Boulder Magazine; BMC Grades (https://thebmc.co.uk/en/indoor-climbing-grades-explained).
Em \textbf{1947}, Fred Bernuz pintou em Fontainebleau o primeiro circuito numerado por cores (\textit{Amarelo, Laranja, Azul, Vermelho, Preto e Branco}), originando os circuitos comerciais modernos \cite{climbski_font, arkose_niveaux}. Paralelamente, consolidou-se a escala alfanumérica \textit{Font} com letras maiúsculas a partir do sexto grau (\texttt{6A} a \texttt{9A}) para distingui-la das vias de corda francesas (\texttt{6a}, \texttt{7b+}) \cite{bmc_grades, actugrimpe_niveaux}.

\subsection*{A.3 John Gill (Escala B) e John Sherman (Escala V)}
% SOURCE [VERIFIED_PRIMARY]: Sherman (1991), Hueco Tanks Guide; Movement Gyms (https://blog.movementgyms.com/understanding-v-grades).
Nos anos 1950, John Gill introduziu o magnésio e a \textbf{Escala B} (\textit{B1, B2, B3}) \cite{movement_vgrades}. Em \textbf{1991}, no guia de Hueco Tanks, John ``Vermin'' Sherman sistematizou a \textbf{Escala Hueco ($V$-Scale)} a partir de $V0$ (posteriormente expandida de \textbf{VB} até \textbf{V17}), hoje padrão nas Américas, África do Sul e Oceania \cite{sherman_hueco, cityrock_guide}.

\begin{table}[!htbp]
\centering
\caption{Síntese histórico-geográfica dos principais sistemas de graduação.}
\label{tab:sistemas_mundo}
\footnotesize
\renewcommand{\arraystretch}{1.12}
\begin{tabularx}{\linewidth}{p{2.2cm}p{2.3cm}p{2.4cm}X}
\toprule
\textbf{Sistema} & \textbf{Origem / Ano} & \textbf{Notação Típica} & \textbf{Escopo e Lógica Estrutural} \\
\midrule
\textbf{UIAA} & Alpes / 1926--1947 & \texttt{V+, VI, VII-, IX+} & Escala em romanos; aberta acima do Grau VI em 1978 \cite{uiaa_history}. \\
\rowcolor{rowgray}
\textbf{Fontainebleau} & França / anos 1940 & \texttt{4, 5+, 6A+, 7B, 8C} & Escala aberta de boulder; letras maiúsculas (\texttt{A, B, C}) a partir do 6 \cite{bmc_grades}. \\
\textbf{YDS} & EUA / 1935--anos 50 & \texttt{5.9, 5.10c, 5.12a} & Sistema decimal de Yosemite; subdividido em \texttt{a--d} a partir de \texttt{5.10} \cite{sportrock_grades}. \\
\rowcolor{rowgray}
\textbf{Ewbank} & Austrália / África do Sul (1967) & \texttt{14, 18, 22, 28, 35} & Escala aberta estritamente numérica contínua para vias \cite{cityrock_guide}. \\
\textbf{Hueco ($V$)} & Texas, EUA / 1991 & \texttt{VB, V0, V3, V8, V17} & Escala aberta padrão para boulder nas Américas \cite{sherman_hueco}. \\
\rowcolor{rowgray}
\textbf{CBME} & Brasil / 2007--2016 & \texttt{5º VIIa E2} / \texttt{VIIb} & Norma DT-2016/01: Grau geral arábico + Crux romano \cite{cbme_dt01}. \\
\bottomrule
\end{tabularx}
\end{table}

\subsection*{A.4 Polos Geográficos Balizadores: Cocalzinho, Igatu e Rocklands}
% SOURCE [VERIFIED_PRIMARY]: Guia de Igatu (EENe, 2014); Climber Codex Pireneus; Guia OuroBoulder (2019); CityROCK (2020).
\begin{itemize}
    \item \textbf{Cocalzinho de Goiás (GO)}: Quartzito compacto de regletes definidos e \textit{mantles} técnicos; referência nacional contra inflação de graus indoor \cite{guia_cocalzinho, guia_ouroboulder}.
    \item \textbf{Vila de Igatu (BA)}: Quartzito áspero de formação tridimensional na Chapada Diamantina, exigindo compressão e pinças em tetos \cite{guia_igatu}.
    \item \textbf{Rocklands (África do Sul)}: Referência global em arenito quartzítico para validação da Escala $V$, espelhada nos ginásios da Cidade do Cabo \cite{cityrock_guide, goodbeta_sa}.
\end{itemize}
